\documentclass{article}
\usepackage[utf8]{inputenc}
\usepackage{amsmath}
\usepackage{amsfonts}
\usepackage{amssymb}
\usepackage{geometry}
\geometry{a4paper, margin=1in}

\title{核心可解数学模型}
\author{Paper 4}
\date{\today}

\usepackage[UTF8]{ctex}
\begin{document}

\maketitle

\section*{保留的核心模型}
根据\texttt{}公式以及PDF文件，可以直接求解的紧凑模型是时间索引混合整数形式（TI），在e-companion中标记为EC.2a--EC.2p。第9页的基于模式的公式、第10页的定价问题以及第22页的连通性切割都属于分支价格与切割框架，但由于它们依赖于隐式路径生成，而不是显式的独立公式化，因此在此处省略了这些内容。

\paragraph{集合与参数。}
设$V=\{0\}\cup N \cup \{n'\}$为节点集，其中$0$为起始仓库，$n'$为虚拟终点仓库。客户集被划分为取货客户$N^{+}$和送货客户$N^{-}$，同时还有$N=N^{+}\cup N^{-}$。设$A$为弧集，$K$为车辆集。对于每一个节点$i$，参数$q_i$表示收集的数量，如果$q_i>0$为真，则表示收集的数量；如果$q_i<0$为真，则表示交付的数量。参数$d_{ij}$表示路由成本，$t_{ij}$表示行驶时间，$Q$表示车辆容量，而$T$表示最大路由持续时间。对于任意子集$S \subseteq N$，设$A(S)=\{(i,j)\in A : i \in S,\ j \in S\}$。

\paragraph{决策变量。}
$x_{ij}^{k}\in\{0,1\}$等于1，当且仅当车辆$k$经过弧线$(i,j)$时。  
$y_i^k$是车辆$k$在客户$i$处收集或交付的数量。  
$z_i^k$是车辆访问节点$i$后的载货量。

\section*{时间索引型多智能体规划公式化}
\begin{equation}
\min \sum_{k \in K} \sum_{(i,j) \in A} d_{ij} x_{ij}^{k}
\tag{EC.2a}
\end{equation}

\begin{equation}
\sum_{(0,i) \in A} x_{0i}^{k} = \sum_{(i,n') \in A} x_{in'}^{k}
\quad \forall k \in K
\tag{EC.2b}
\end{equation}

\begin{equation}
\sum_{(i,j) \in A} x_{ij}^{k} = \sum_{(j,i) \in A} x_{ji}^{k}
\quad \forall i \in V,\ \forall k \in K
\tag{EC.2c}
\end{equation}

\begin{equation}
\sum_{k \in K} y_i^k = q_i
\quad \forall i \in V
\tag{EC.2d}
\end{equation}

\begin{equation}
y_i^k \le \min\{q_i,Q\} \sum_{(i,j) \in A} x_{ij}^{k}
\quad \forall i \in N^{+},\ \forall k \in K
\tag{EC.2e}
\end{equation}

\begin{equation}
y_i^k \ge \max\{q_i,-Q\} \sum_{(i,j) \in A} x_{ij}^{k}
\quad \forall i \in N^{-},\ \forall k \in K
\tag{EC.2f}
\end{equation}

\begin{equation}
z_j^k \ge z_i^k + y_j^k + (x_{ij}^{k} - 1)(Q+1)
\quad \forall j \in N^{+},\ \forall (i,j) \in A,\ \forall k \in K
\tag{EC.2g}
\end{equation}

\begin{equation}
z_j^k \le z_i^k + y_j^k + (1 - x_{ij}^{k})(Q+1)
\quad \forall j \in N^{-},\ \forall (i,j) \in A,\ \forall k \in K
\tag{EC.2h}
\end{equation}

\begin{equation}
z_i^k \le Q
\quad \forall i \in V,\ \forall k \in K
\tag{EC.2i}
\end{equation}

\begin{equation}
\sum_{(i,j) \in A} t_{ij} x_{ij}^{k} \le T
\quad \forall k \in K
\tag{EC.2j}
\end{equation}

\begin{equation}
\sum_{(i,j) \in A(S)} x_{ij}^{k} \le |S| - 1
\quad \forall S \subseteq N,\ |S| > 1,\ \forall k \in K
\tag{EC.2k}
\end{equation}

\begin{equation}
y_0^k = y_{n'}^k = 0
\quad \forall k \in K
\tag{EC.2l}
\end{equation}

\begin{equation}
x_{ij}^{k} \in \{0,1\}
\quad \forall (i,j) \in A,\ \forall k \in K
\tag{EC.2m}
\end{equation}

\begin{equation}
y_i^k \ge 0
\quad \forall i \in N^{+},\ \forall k \in K
\tag{EC.2n}
\end{equation}

\begin{equation}
y_i^k \le 0
\quad \forall i \in N^{-},\ \forall k \in K
\tag{EC.2o}
\end{equation}

\begin{equation}
z_i^k \ge 0
\quad \forall i \in V,\ \forall k \in K
\tag{EC.2p}
\end{equation}

\paragraph{当前问题集网络惯例。}
TI的公式中使用了一种封闭的有向网络。方程(EC.2c)适用于$V$中的每个节点，包括两个仓库的副本。所提供的实例中使用了$A=\{(i,j):i,j\in V,\ i\ne j\}$，其中包括$(n',0)$以及其他仓库之间的连接线。方程(EC.2a)和(EC.2j)包含了所有选定连接线的成本和时间，包括返回路径。客户数量满足$q_i\ge0$对于$i\in N^+$，以及$q_i\le0$对于$i\in N^-$的条件，这意味着零需求是可能的，同时$q_0=q_{n'}=0$的条件也得到满足。 

\end{document}
